\documentclass{situationist-preprint}

\title{Hume Looked in the Wrong Place}
\msid{PID-04-b}
\msfamily{PID-04}
\msdiscipline{Personal Identity and the Self}
\mstype{Companion manuscript}
\mskeywords{Hume; bundle theory; introspection; self; organization; category mistake}

\begin{document}
\maketitle

\begin{abstract}
Hume reported that when he entered most intimately into what he called himself, he found only particular perceptions and never a self. The report is commonly taken as evidence that there is no self, or none beyond a bundle. PID-04-a holds that the self is real, a construction: the organization of self-representation, self-interpretation and self-regulation by which an organism takes itself as a single subject and agent. We argue that Hume's search could not have found the self so understood, because it looked for an item among the perceptions, and the self is their organization. The negative result of the search is therefore no evidence against the self's existence, and Hume's own later doubts about the bundle view point to the same conclusion.
\end{abstract}

\section{The search}

\citet[1.4.6]{Hume1739} reasoned as follows. Every real idea is derived from some impression. If any impression gave rise to the idea of self, it would have to remain invariably the same through the whole course of our lives. No impression is constant and invariable. So there is no such idea, and the self is nothing but a bundle of different perceptions in perpetual flux. The decisive step is the report of introspection: whenever he looked, he stumbled on a perception of heat, cold, love or hatred, and never caught himself without a perception, nor observed anything but the perception.

The search presupposes a view of what it was looking for. It treats the self as a candidate object of inner observation, an item that, if it existed, would be found alongside the perceptions, distinguished from them by its constancy. Call this the \emph{item assumption}.

\section{Why the search had to fail}

\begin{proposition}\label{p:fail}
If the self is as PID-04-a, Definition~2.3, describes it, no act of introspection presents the self as one of its objects alongside the perceptions, and none could.
\end{proposition}

\begin{proof}
On Definition~2.3 of PID-04-a, the self is the unified, ongoing structure by which an organism represents itself as a single subject, ascribes its states to that subject and governs its conduct by those ascriptions. A structure of this kind is the way the perceptions and actions are organized: which states are ascribed to whom, how they are integrated, which guide which. It is not a further state with a content of its own. An act of introspection presents states with contents. So the self is not among the things introspection presents. Moreover, the act of introspection is itself one of the activities the structure organizes, since in introspecting the subject ascribes the perceptions it finds to itself. It cannot present, as an object, the organization within which it occurs.
\end{proof}

The point is \citeauthor{Ryle1949}'s, applied inward. A visitor shown the colleges, libraries and laboratories who then asks where the university is has made a category mistake: the university is the way the colleges and libraries are organized, not a further building \citep[ch.~1]{Ryle1949}. Hume, looking among the perceptions for the self, made the same mistake. Looking among the perceptions for the self is like looking among the players for the team.

The point also explains why the search seemed so conclusive. Every time Hume looked, he found some perception, and he found it as his: the search itself depended on his ascribing what he found to the one who was searching. What cannot be bracketed in such an act of reflection is the one who performs it, present in each act and not as an object of it (PHE-20, PID-13). This presence is not a persisting item the search overlooked. It is the organization at work in the search.

\section{Hume's second thoughts}

In the Appendix to the \emph{Treatise}, Hume confessed himself dissatisfied with his account. Having loosened the perceptions into distinct existences, he could not explain what principle unites them into one mind, and he found the difficulty too hard for his understanding \citep[Appendix]{Hume1739}. The difficulty is a symptom of the item assumption. If the self must be an item and no item is found, all that is left are distinct perceptions, and their unity becomes inexplicable. If the self is the organization of the perceptions, their unity is just what the self consists in, and there is nothing left to explain by a further principle. Hume's own comparison of the mind to a republic, whose members are united by relations of government and subordination and which retains its identity while its members change \citep[1.4.6]{Hume1739}, describes an organization, not a bundle.

\begin{objection}
If the self is merely the organization of perceptions, Hume was right after all: the self is nothing over and above the bundle.
\end{objection}

\begin{reply}
An organization is not nothing over and above its members. A bundle is a mere collection; an organization is a structure that figures in explanations the members alone do not supply. By PID-04-a, Lemma~3.2, the self figures in true explanations that cannot be reformulated without it, so by Principle~P1 it exists. What is right in Hume is that the self is not a simple, unchanging substance; what is wrong is the inference from its not being an item to its not being anything.
\end{reply}

\section{Conclusion}

Hume's search was guided by the assumption that the self, if real, would be an object of inner perception. On the view of PID-04-a it is the organization of perception and action under self-ascription, and an organization is not perceived among the things it organizes. The search failed because it looked in the wrong place; its failure tells against the item assumption, not against the self.

\begin{thebibliography}{99}\small
\bibitem[Hume(1739)]{Hume1739} Hume, D. (1739). \emph{A Treatise of Human Nature}. Cited by book, part and section.
\bibitem[Ryle(1949)]{Ryle1949} Ryle, G. (1949). \emph{The Concept of Mind}. Hutchinson.
\end{thebibliography}

\end{document}
